\chapter[Homeostasis: Constant Correction, Not Constancy]{Homeostasis:\protect\\ Constant Correction, Not Constancy}

\begin{kaichang}
Put a thermometer under your arm for five minutes. Whether it is a summer afternoon or winter morning, whether you just drank hot soup or left an air-conditioned room, the reading will probably fall in the narrow interval between $36\,^{\circ}\mathrm{C}$ and $37\,^{\circ}\mathrm{C}$.

Think how strange that is. Outdoor temperatures can vary by $40\,^{\circ}\mathrm{C}$ over a year. Food is hot or cold, and you alternate sprinting a hundred meters with lying asleep. Your internal boiler's output keeps changing, yet the thermometer barely moves. Is your body really holding temperature still?

Guess an answer. By the end, you will discover the opposite: temperature, glucose, pH, and water content deviate and fluctuate every moment. An endlessly working correction team pulls them back. Stability means not immobility, but promptly correcting movement.
\end{kaichang}

\section{Those Unruly Numbers}

Start with measurements. A thermometer, glucose meter, blood-pressure monitor, and laboratory report reveal a series of numbers:

\begin{itemize}[itemsep=2pt]
\item Temperature: roughly $36.1\,^{\circ}\mathrm{C}$ to $37.2\,^{\circ}\mathrm{C}$, fluctuating by less than $1\,^{\circ}\mathrm{C}$ daily, lower in the morning and higher in the evening.
\item Fasting glucose: about $0.7$ to $1.0$ grams per liter of blood, rising after meals and returning within two hours.
\item Blood pH: $7.35$ to $7.45$, an alarmingly narrow range; a deviation of $0.4$ can endanger life.
\item Blood sodium: about $135$ to $145$ millimoles per liter, around which thirst and increased urination revolve.
\item Oxygen saturation: usually above $95\,\%$ in healthy people at low altitude.
\end{itemize}

Notice two things. These are \textbf{ranges}, not fixed points. And you challenge them daily: a sugary drink raises glucose, an eight-hundred-meter run makes muscles produce lactate and heat, and hot pot delivers plenty of salt. Yet hours later, the numbers return.

Think of a tightrope walker. From afar, the person seems stable; close up, body and balance pole constantly adjust from side to side. Your body is that walker, balancing on five or six ropes simultaneously for a lifetime.

\section{Zooming In: What Surrounds Cells?}

Why such vigilant correction? Return to particles.

Chapter 57 explained that your tens of trillions of cells are not solid bricks packed together, but bathed in thin layers of fluid: blood and tissue fluid, together the \textbf{internal environment}. This fluid delivers each cell's supplies: glucose, oxygen, and ions are taken from it, while carbon dioxide and wastes enter it.

Inside is a precise chemical factory. Chapters 41--43's enzymes are fussy proteins: a few degrees of warming loosens their shapes and sharply lowers activity; more heat denatures them irreversibly like cooked egg white in Chapter 44. A pH shift changes surface charges and prevents substrates binding. Sodium, potassium, and calcium concentrations are vital to nerves and muscles: Chapter 58's nerve impulses are waves of ions crossing membranes. Wrong concentrations can disrupt heartbeat.

\textbf{The cell's reaction network operates under demanding environmental specifications}. Departures slow the factory, cause errors, or halt production. Coma at a fever of $41\,^{\circ}\mathrm{C}$ reflects countless enzymes failing together, not a body being too delicate.

These numbers must be watched. Who watches, and how are they corrected?

\section{The General Blueprint: Negative Feedback}

Nature, or evolution in the chapters from 83 onward, repeatedly uses \textbf{negative feedback}. It has three roles, familiar from your air conditioner:

\begin{itemize}[itemsep=2pt]
\item \textbf{Sensor:} the air conditioner's temperature probe; bodily temperature and chemical receptors report current values.
\item \textbf{Control center:} its circuit board; your hypothalamus, a fingernail-sized brain region, pancreatic islet cells, and others compare readings with a \textbf{set point} and determine the error.
\item \textbf{Effector:} its compressor; sweat glands, muscles, liver, kidneys, and vessels pull values back.
\end{itemize}

Negative means correction opposes deviation: push high values down and lift low ones up. Suppose summer sunshine raises your temperature by $0.5\,^{\circ}\mathrm{C}$. Skin and hypothalamic receptors sense warming; the hypothalamus orders sweating and skin-vessel dilation, carrying heat to the surface for evaporation. Temperature falls, and the command ends. Winter reverses the response: skin vessels constrict to reduce heat loss and muscles shiver involuntarily. Shivering is skeletal muscle doing work without useful movement, turning ATP energy into heat as temporary human heating.

Glucose uses the same plan with different participants. After a meal, islet $\beta$ cells act as combined sensor and controller, releasing insulin to urge liver, muscle, and fat cells to store glucose, as Chapter 55 detailed. During hunger, $\alpha$ cells release glucagon, ordering liver glycogen breakdown and glucose release. Opposing teams alternate duty and confine glucose to a narrow band.

\begin{qiaoliang}
How hard does correction work? With about $5$ liters of blood and an upper value of $1$ gram of glucose per liter, \textbf{only about $5$ grams of glucose circulate in your vessels}, less than a small sugar cube. A can of cola contains about $35$ grams of sugar, $7$ times that entire supply. Without correction, it could theoretically raise glucose to $8$ times normal. Yet two hours after eating, values largely return, as tens of grams are delivered to liver and muscle stores within tens of minutes. Consider heat too: resting cellular respiration in Chapter 53 produces about $80$ to $100$ watts, like an old incandescent bulb. If none escaped, a $70$-kilogram body with heat capacity near $245$ kilojoules per degree Celsius, and specific heat near water's, would warm about $1.4\,^{\circ}\mathrm{C}$ hourly, reaching danger within half a day. Continuous matching of heat loss to production prevents you from cooking yourself.
\end{qiaoliang}

Negative feedback aims not to freeze at one point but to fluctuate slightly around a set point. Real glucose and temperature traces are wavy, with tightly limited amplitudes. The title's full meaning is \textbf{homeostasis means correcting deviations as they arise, not remaining unchanged}.

\section{Sometimes Reinforcement Helps: Positive Feedback}

Does the reverse exist, reinforcing larger deviations? Yes: \textbf{positive feedback}, used when a process must proceed decisively to completion.

Childbirth is classic. Uterine contractions press the fetus against the cervix, whose nerve signals reach the brain and prompt more oxytocin. Oxytocin strengthens contractions, driving another round until delivery physically ends the cycle. Clotting is similar: activated factors activate more partners, rapidly amplifying like dominoes to plug the opening before excessive blood loss.

Positive feedback is \textbf{local and has an endpoint}. Delivery ends labor; wound healing stops clotting. If it ruled indefinitely like negative feedback, it would drive a variable to extremes: loss of control. Negative feedback is the default; positive feedback is a tightly contained strike team.

\section{Fast Nerves and Slower Hormones}

How are commands delivered? Two contrasting, complementary routes from Chapter 58 divide the work.

\textbf{Nerves} are dedicated point-to-point lines. Electrical impulses travel tens of meters per second, arriving in milliseconds, precisely targeted and quickly withdrawn. Shivering, vessel constriction, and faster heartbeat use this route for immediate responses.

\textbf{Endocrine signaling} broadcasts more slowly. Hormones enter blood and circulate body-wide, arriving in seconds to minutes. Everyone receives them, but only receptor-bearing cells open the letters, Chapter 58's lock-and-key model. Insulin, glucagon, antidiuretic hormone, and thyroid hormone use this route. Its strengths are broad coverage and duration: adjusting glucose, water, and salts requires dozens of cell types, too many for dedicated wiring alone.

Fast responses and lasting control meet in the hypothalamus. Receiving temperature and osmotic-pressure reports, it sends nerve commands and directs the pituitary hormone dispatch center. It is your round-the-clock control room.

\section{The Chemists' Territory: Kidneys, Lungs, and Liver}

Beyond temperature and glucose, the most chemical correction occurs in three organs controlling internal-environment \textbf{composition}.

\textbf{Lungs} regulate two gases. Expelling too much carbon dioxide makes blood more alkaline; too little makes it more acidic. The buffering system discussed next responds first, while lungs provide subsequent adjustment. Faster, deeper breathing during running is not merely breathlessness but quicker removal of respiration's \chem{CO_2}. Chemoreceptors in the neck's carotid arteries monitor oxygen and \chem{CO_2}; falling oxygen, at altitude or during hard exercise, prompts respiratory muscles through nerves to restore intake.

\textbf{The liver} distributes chemicals and detoxifies. It stores and releases glucose in Chapter 55, breaks down surplus amino acids, converts toxic ammonia into less toxic urea for renal disposal, and modifies and inactivates much alcohol and many drugs.

\textbf{Kidneys} tirelessly check quality, filtering the body's blood dozens of times daily to produce about $180$ liters of initial filtrate, more than two large water-cooler bottles. They recover $99\,\%$ of water, all glucose, and needed salts, concentrating surplus water, salts, urea, and other leftovers into one or two liters of urine. A large drink soon makes urine more abundant and dilute; a salty meal makes it concentrated and triggers thirst. Kidneys and thirst jointly hold sodium between $135$ and $145$. Antidiuretic hormone is this loop's slow messenger: concentrated blood prompts hypothalamic release, instructing kidneys to recover more water.

In short, lungs manage gases, liver transformations, and kidneys exchange. Together they maintain the internal fluid $24$ hours a day.

\section{Symbols: Buffering Acid--Base Shocks}

Now express pH regulation symbolically. Blood's leading buffer partners are carbon dioxide and bicarbonate, linked by reversible equilibria:

\[\chem{CO_2} + \chem{H_2O} \rightleftharpoons \chem{H_2CO_3} \rightleftharpoons \chem{H^{+}} + \chem{HCO_3^{-}}\]

Recall Chapter 28's equilibrium-shift rule. This is a chemical spring: a sudden \chem{H^{+}} increase, such as lactic-acid dissociation during exercise, shifts it left. \chem{HCO_3^{-}} captures extra \chem{H^{+}}, forming \chem{CO_2} exhaled by the lungs. Reduced \chem{H^{+}} shifts it right to replenish protons. As long as the partners are not exhausted, pH remains between $7.35$ and $7.45$.

Notice external support at both ends: lungs control \chem{CO_2} removal, while kidneys recover or excrete \chem{HCO_3^{-}}, more slowly over hours. Chemical equilibrium handles the first seconds, lungs adjust over minutes, and kidneys follow over hours: three defenses on different timescales. The general principle is \textbf{fast systems hold the line; slow ones finish the job}.

\section{Evidence: From Internal Environment to Homeostasis}

\begin{zhengju}
Maintaining a stable internal environment sounds obvious today, but in the mid-nineteenth century it was hardly a recognized question. The default view was that blood contents came directly from food and drink: glucose rose simply because sugar had been eaten.

The French physiologist Claude Bernard challenged this. He compared sugar in the livers and blood of normally fed dogs with fasting dogs. If all sugar came directly from food, fasting livers should contain none. Instead, \textbf{fasting livers still released sugar into blood}. Further investigation identified a storage substance, later named glycogen, that the liver converted to glucose. Blood sugar was not just food's reflection: the body \textbf{produced} it, adjusting production to need.

Bernard proposed a larger idea: a multicellular organism's true living environment is not outside weather but its internal fluid, the milieu intérieur. In his 1865 lectures, he wrote that constancy of this environment permits free, independent life. This was no philosophical musing but a conclusion from measurements such as sugar in fasting dogs' livers.

Half a century later, the American physiologist Cannon asked how constancy was achieved. He and students stimulated animal sympathetic nerves and measured linked changes in glucose, heartbeat, and adrenal secretion, finding responses that opposed deviations. In 1926, he named the mechanism homeostasis, Greek for similar plus standing, deliberately not identical: values were held within ranges rather than frozen. Later, control engineers recognized the same negative-feedback blueprint in thermostats and cruise control. Physiology and engineering independently drawing the same loop suggests a genuine general rule.
\end{zhengju}

\section{Fever: A Raised Set Point, Not a Runaway Boiler}

Homeostasis offers a surprising view of fever.

A fever of $39\,^{\circ}\mathrm{C}$ \textbf{does not} mean failed cooling. Heating and cooling still work, but the maintained set point has been deliberately raised. Your sensations provide evidence: as fever starts, you feel \textbf{cold} and shiver under blankets despite rising temperature. The hypothalamus has already shifted its target to $39\,^{\circ}\mathrm{C}$ while the body is only at $38\,^{\circ}\mathrm{C}$. Relative to the new target, you are cold, triggering shivering, constricted vessels, and a desire for blankets. When the target returns to $37\,^{\circ}\mathrm{C}$, actual temperature is now too high, triggering heavy sweating. Chills then heat reflect a moving set point followed by correction.

Who changes it? Immune-cell signals, cytokines from Chapter 61, reach the hypothalamus and direct local prostaglandin synthesis, turning up the thermostat. Why has evolution retained this? Evidence largely suggests moderate warming slows some pathogens and improves immune-cell activity. Fever resembles deliberately raising workshop temperature to fight infection rather than an accident.

This does not mean every fever should be ignored. Very high temperatures, such as above $40\,^{\circ}\mathrm{C}$, prolonged fever, or fever in babies, older people, and those with underlying illness can make risks outweigh benefits. Fever-reducing drugs lower the set point. Medication and care decisions depend on the person and circumstances, not one fixed number. Understanding the mechanism matters more than memorizing a drug threshold.

\textbf{Translator safety note:} Do not wait for $40\,^{\circ}\mathrm{C}$ before seeking help. A baby under 3 months with $38\,^{\circ}\mathrm{C}$ or higher needs urgent medical advice; other concerning symptoms also matter. Follow \href{https://www.nhs.uk/symptoms/fever-in-children/}{NHS child-fever guidance}, not the simplified mechanism as a treatment rule.

\section{Where the Model Breaks Down}

Homeostasis is useful enough that its boundaries must be clear.

\begin{itemize}[itemsep=2pt]
\item \textbf{Correction has limits.} In hot, humid conditions, sweat cannot evaporate and cooling is physically blocked, so heat production drives temperature upward. Heat illness means cooling correction has lost: the body cannot adjust, rather than choosing not to. Prolonged extreme cold and severe diarrheal dehydration are similar.
\item \textbf{Parts fail.} Destroyed $\beta$ cells or insulin insensitivity impair the glucose loop's controller, leaving glucose chronically high. Diabetes is a broken negative-feedback loop. Treatment supplies missing insulin instructions or improves sensitivity or glucose lowering through medication.
\item \textbf{Set points are not universal constants.} Normal ranges come from population statistics. Someone may normally be at $36.0\,^{\circ}\mathrm{C}$ and another at $36.8\,^{\circ}\mathrm{C}$. Aging weakens temperature regulation and thirst, increasing heat-wave vulnerability; immature infant systems fluctuate more. Women's temperatures also vary with the menstrual cycle.
\item \textbf{The stable range can change.} Long-term high-altitude residents increase red-cell counts to adapt to low oxygen, moving the production line rather than restoring lowland oxygen values. Allostasis describes such longer-term resetting of operating points. Set points themselves arise through evolution and adaptation, not immutable inscriptions.
\end{itemize}

\textbf{Translator safety note:} Heatstroke is not an ordinary fever. Confusion or loss of consciousness during heat exposure requires emergency help and prompt cooling; call your local emergency number rather than waiting for feedback or fever medicine to work. See \href{https://www.cdc.gov/niosh/heat-stress/about/illnesses.html}{CDC heat-illness guidance}.

A further limit: we separated temperature, glucose, and pH loops for explanation. Real loops intertwine: low glucose affects temperature, dehydration affects blood pressure and kidneys, and fever changes heartbeat and breathing. Separate diagrams do not mean separate bodily offices.

\begin{wuqu}
``Sweating detoxifies; more sweat cleans the body.'' Sweat is over $99\,\%$ water with small amounts of sodium, potassium, and urea. Its main job is cooling: each evaporated gram removes about $2.4$ kilojoules. Metabolic wastes, drugs, and alcohol mainly follow liver transformation and kidney excretion; sweat's small urea concentration is insignificant beside urine's. Deliberately bundling up or forcing heavy sweating in heat does not detoxify but loses water and salts, precisely what the internal environment regulates. You \textbf{create} new problems for correction. Drinking appropriately and allowing normal liver and kidney function does more than forced sweating.
\end{wuqu}

\section{Try It: Watch Negative Feedback}

\begin{shiyan}
\textbf{Materials:} a phone stopwatch and a quiet seat.

\textbf{Procedure:}
\begin{enumerate}[itemsep=1pt]
\item Sit quietly for $5$ minutes. Feel the radial pulse on the thumb side of your inner wrist, count for $30$ seconds, multiply by $2$, and record resting heart rate.
\item Do $1$ minute of moderate exercise, such as squats or jumping jacks, then immediately count the pulse for $30$ seconds and record post-exercise heart rate.
\item Measure every $1$ minute for another $4$ to $5$ readings and connect them on a line graph.
\end{enumerate}

\textbf{Observations:} how high did the rate rise, and how long to return near rest? Does it drop abruptly or approach gradually? Who commanded the increase, and what deviation was addressed, namely increased muscle oxygen demand? Which feedback stage corresponds to recovery? What could cause slow recovery?

\textbf{Safety:} remain slightly breathless but able to speak. Stop and rest immediately for dizziness or chest tightness. Anyone with heart disease or medical advice against exercise should measure only resting pulse and analyze curves, without forcing activity.
\end{shiyan}

\section{Back to the Thermometer}

That steady thirty-six-point-something reading now reports the correction team's work. During the five-minute measurement, sweat glands adjusted output, skin vessels adjusted diameter, the liver quietly released glucose from glycogen, kidneys decided water recovery drop by drop, and the hypothalamus monitored a dozen signals like a control tower. The thermometer's stillness comes from countless moving molecules.

This also explains why laboratory reports list reference ranges instead of single standards. Medical testing fundamentally \textbf{samples the internal environment}: glucose, sodium, uric acid, and transaminases ask whether correction still holds. A slightly out-of-range value need not cause immediate alarm, since ranges reflect populations, individual variation, and measurement error. But persistent, substantial deviations often indicate a sensor, controller, or effector problem worth investigating medically. Relating each result to a loop can make reports less mysterious.

\section{A Special Branch of the Correction Team}

So far, correction has concerned numerical variables: temperature, concentration, and pH. Other threats are \textbf{living invaders}, bacteria, viruses, and parasites, which replicate, mutate, and evade rather than follow your negative feedback. Restoring a set point is insufficient. You need a force that \textbf{distinguishes friend from foe, removes targets, and remembers enemies}.

That is next chapter's subject. Fever-raising signals are its messengers, and vaccines work through its memory. Homeostasis regulates internal physical chemistry; immunity maintains security. The systems continually communicate through blood and hypothalamus. Chapter 61 introduces this learning army.

\begin{tiaozhan}
One algebraic line reveals negative feedback. Let error from a set point be $E$, let each round remove fraction $k$ with $0<k<1$, and let the environment add disturbance $D$ each round. Then $E' = (1-k)E + D$. At steady state, $E = (1-k)E + D$, giving $E = D/k$. Three lessons follow. First, persistent $D$ means error \textbf{never reaches zero}: steady state is dynamic balance with a small error, not perfection. Second, larger $k$ means smaller residual error and flatter glucose or temperature regulation. Third, as $k$ approaches $0$ through control failure, error explodes; diabetes is a collapse of $k$ with divergent $D/k$. Real physiology is more complex, with nonlinear correction and delays. Delay plus excessive correction produces oscillation, explaining swinging glucose or temperature in some patients. This box is optional, but next time you hear dynamic balance, picture $E=D/k$.
\end{tiaozhan}

\section*{Practice}

\begin{sikaoti}
\item Draw sensor → controller → effector → sensor for sodium regulation after drinking a large bottle of water. Identify each participant and the signals connecting them.
\item Plot set-point and actual-temperature curves over fever onset, plateau, and resolution. Explain initial chills and later sweating as the set point moves first and temperature follows.
\item During strenuous exercise, muscles produce heat, \chem{CO_2}, and lactate. List at least three simultaneous correction loops for heat, \chem{CO_2}, and \chem{H^{+}}, with effectors and approximate timescales.
\item Use homeostasis as small fluctuations around $E=D/k$, not immobility, to refute the claim that healthy glucose must be a straight line and any fluctuation indicates disease. How large a fluctuation warrants concern, and what evidence informs your judgment?
\item Positive feedback becomes disastrous without termination. Identify the stop mechanisms in childbirth and clotting. What would uncontrolled clotting do inside vessels?
\item Look up three values and reference ranges from your or a family member's latest medical report. For each, describe the likely correction loop, sensor, controller, and effector.
\end{sikaoti}
